\documentclass[8pt,xcolor=dvipsnames]{beamer}
\usepackage[latin1]{inputenc}

\usetheme[height=7mm]{Berlin}
\usecolortheme[named=Brown]{structure} 


\usepackage{time}       % date and time
\usepackage{graphicx}
\usepackage[T1]{fontenc}    % european characters
%\usepackage{courier}
\usepackage{animate}
\usepackage{multirow}
\usepackage{natbib}
\usepackage{amssymb,amsmath}  % use mathematical symbols
\usepackage{bookman}                  % use palatino as the default font
\setbeamercovered{transparent}
%\newcommand{\figpath}{/home/phybva/SiOJets_New/TALK/NEWFIGS}
\newcommand{\figpath}{./NEWFIGS/}
\newcommand{\spic}[2]{\includegraphics[scale=#1]{#2}}
\newcommand{\myref}[1]{{\small{\color{red}{(#1)}}}}
\newcommand{\movauto}[5]{\animategraphics[width=#1,height=#2,autoplay]{#3}{#4}{0000}{#5}}
\newcommand{\movloop}[5]{\animategraphics[width=#1,height=#2,loop]{#3}{#4}{0000}{#5}}

\title[YSO Outflows]{Outflows from Young Stellar Objects}
\subtitle{MHD, Radiation \& Chemistry}
\author[Bhargav Vaidya]{\textcolor{blue}{Bhargav Vaidya}\inst{1}}
\institute[Uni. Leeds]{\inst{1}School of Physics and Astronomy, University of Leeds, Leeds.}
\date[2014]{{\large{Astro Seminar, TIFR}}\\ {\small{January 13, 2014}}}
%\titlegraphic{\vskip30pt \hskip215pt\includegraphics[width=0.3\textwidth]{\figpath/Leeds-logo.jpg}

\begin{document}

\begin{frame}
\vspace*{0.5cm}\titlepage  
  \vspace*{0.1cm}\textcolor{red}{{\textit{Collaborators:}}}\\
  \tiny{Paola Caselli (Uni. Leeds), Thomas Douglas (Uni. Leeds), Oliver Porth (Uni. Leeds), \\
  Christian Fendt (MPIA), Henrik Beuther (MPIA), Somayeah Sheiknezami (MPIA),\\
  Ciraco Goddi (JIVE),  Andrea Mignone (Uni. Torino).\\}
%\titlepage
\end{frame}

\begin{frame}
\frametitle{Outline}
\tableofcontents
\end{frame}

\section{Introduction}
\begin{frame}{Star formation : What do we know}

\end{frame}

\begin{frame}{Star formation : Feedback}
Talk about winds and outflows
\end{frame}


\begin{frame}{Present challenges}
\end{frame}


\section{Motivation}
\begin{frame}{Outflow Evolutionary Picture}
\end{frame}

\begin{frame}{Chemistry in outflows}
\end{frame}

\begin{frame}{Molecular bullets and EHV emission}
\end{frame}

\section{Methods : Numerical Simulations}
\begin{frame}{Numerical code}

\end{frame}

\begin{frame}{Launching and Propagation}
\end{frame}

\begin{frame}{Radiation force}
\end{frame}

\begin{frame}{Chemistry and Cooling}
\end{frame}

\section{Outflow Dynamics : Launching}
\begin{frame}{MHD Acceleration}
\end{frame}

\begin{frame}{Collimation and Radiation}
\end{frame}

\begin{frame}{Force Parameters and its impact}
\end{frame}

\begin{frame}{Resistive effects}
\end{frame}

\begin{frame}{A case of Orion Source I}
\end{frame}

\section{Outflow dynamics : Propagation}

\begin{frame}{Cooling in Jets}
\begin{columns}
\begin{column}{0.3\textwidth}
\vspace{-3cm}
\begin{block}{Cooling is important}
\begin{itemize}
\item t$_{\rm cool} \sim n_0/\Lambda(T) <$ t$_{\rm dyn}$
  \myref{e.g, Blondin, 1990}\\
\item Cooling rate for a molecular medium with n$_0$ = 10$^{5}$ cm$^{-3}$
\end{itemize}
\end{block}
\end{column}
\begin{column}{0.6\textwidth}
\spic{0.3}{\figpath/pfig1.pdf}
\end{column}
\end{columns}
\end{frame}

\begin{frame}{Cooling in Jets}
\begin{block}{Dynamical Effects}
\centering
\begin{itemize}
\item Thinner jets with less pronounced coccoon
\item Enhanced density in the internal knots with instable features.
\item Thermal Instability in bow-shock of the jet. 
\end{itemize}
\end{block}
\centering
\spic{0.23}{\figpath/pfig4.pdf}

\end{frame}

\begin{frame}{Molecular Interplay}
\begin{columns}
\begin{column}{0.4\textwidth}
\spic{0.2}{\figpath/pfig6.pdf}
\end{column}
\hfill
\begin{column}{0.6\textwidth}
\spic{0.2}{\figpath/pfig7.pdf}
\end{column}
\end{columns}
\end{frame}



\begin{frame}{SiO Abundance and Jet Velocity}
\end{frame}

\begin{frame}{Spectra and PV diagrams : SiO (2-1)}
\begin{columns}
\begin{column}{0.48\textwidth}
\spic{0.17}{\figpath/pfig8.pdf}\\
\hspace{0.2cm} $\phi = \pi/2$ (Plane of Sky)
\end{column}
\hfill
\begin{column}{0.48\textwidth}
\spic{0.17}{\figpath/pfig9.pdf}\\
\hspace{0.2cm}$\phi = \pi/4$ 
\end{column}
\end{columns}
\end{frame}

\begin{frame}{Multi-Line survey : Emission I}
\small{
\begin{itemize}
\item Molecular Jet with $\eta$ = 3.0\\
\item Line Intensities with Top Hat Profile. \\
\item Jet is in plane of sky and convolved with a Gaussian beam of 2.5''
\end{itemize}
}
\begin{tabular}{ccc}
\textbf{J = 2-$>$1} & \textbf{J = 5-$>$ 4} & \textbf{J = 8-$>$ 7}\\
\movauto{3.5cm}{5.5cm}{3}{Img21/Img21_}{0010}
&
\movauto{3.5cm}{5.5cm}{3}{Img54/Img54_}{0010}
&
\movauto{3.5cm}{5.5cm}{3}{Img87/Img87_}{0010}

\end{tabular}

\end{frame}

\begin{frame}{Multi-Line survey : Emission II}
\end{frame}

\begin{frame}{Case of L1448 \myref{Nisini 2007}}
\begin{columns}[t]
\begin{column}{0.25\textwidth}
\spic{0.1}{\figpath/Nisini2007_f1.pdf}\\
Contours SiO 2-1
\end{column}
%\hfill

\begin{column}{0.40\textwidth}
 \spic{0.25}{\figpath/Nisini2007_f4.pdf}\\
 Spectral Features
\end{column}
%\hfill

\begin{column}{0.25\textwidth}
\spic{0.25}{\figpath/Nisini2007_f6.pdf} \\
Kinematic Study (LVG) 
\end{column}
%\hfill

\end{columns}
\end{frame}

\begin{frame}{Multi-line survey : Line Ratios}

\begin{columns}[T]
\begin{column}{0.35\textwidth}
\begin{itemize}
\item EHV emission of 0.5 K.
\item Line ratios close to Unity. 
\item Multi-line emission show a distinct fall at high J$_{\rm up}$.
\end{itemize}
\end{column}
\hfill


\begin{column}{0.75\textwidth}
\spic{0.28}{\figpath/pfig14.pdf} 
\end{column}
\hfill


\end{columns}

\end{frame}

\begin{frame}{Predictions for ALMA}
\begin{block}{ALMA Cycle 2}
\textbf{Bands 3, 6, 7}: Assuming a source placed at 400 pc.[Orion - HH 212 \myref{Zinnecker,1998}] 
\end{block}
\begin{center}
\spic{0.21}{\figpath/pfig12.pdf}
\end{center}
\end{frame}

\begin{frame}{Focussing on a single knot}
\begin{tabular}{c c}
\spic{0.2}{\figpath/codella_hh212_f3.pdf} & \spic{0.2}{\figpath/knot_ang85.pdf} \\
HH 212 \myref{Codella 2007} &\\
\end{tabular}
\end{frame}

\begin{frame}[t]{Rotation or Wiggles?}
  \begin{center}
  \movloop{8cm}{5cm}{61}{Img_3djet/contourjet1.}{0060}
  \end{center}
\end{frame}



\section{Summary}
\begin{frame}{Conclusions}
\end{frame}

\end{document}

